\documentclass[conference]{IEEEtran}
\usepackage[T1]{fontenc}
\usepackage{amsmath,amssymb}
\usepackage{booktabs}
\usepackage{graphicx}
\usepackage[hidelinks]{hyperref}

\title{Copilot Interaction Mode and Productivity Across Programmer Experience Levels}
\author{\IEEEauthorblockN{Damir Makhmutov and Amir Kasimov}
\IEEEauthorblockA{Research Methods Course Project}}

\begin{document}
\maketitle
\begin{abstract}
AI-assisted programming raises questions about when suggestions should be
presented and whether interaction preferences differ with programmer experience.
This work describes an exploratory secondary analysis of published GitHub
Copilot productivity measurements. The dataset contains 66 task sessions from
33 participants, with one automatic-suggestion session and one manual-invocation
session per participant. Both conditions provide access to Copilot. We define
productivity as completed task fraction per hour and compare automatic-minus-manual
changes within novice and expert groups. Our primary estimand is the difference
between the groups' mean paired changes. The analysis validates source checksums,
recalculates productivity, preserves participant pairs, and uses stratified
participant-level bootstrap intervals. Completion and time are reported separately
to make the productivity ratio interpretable. This laboratory-stage draft presents
the introduction and implemented methodology; final results, scientific figures,
and sensitivity analyses will be integrated after Task 3. It makes no claim of a
new participant experiment, an AI-free control, or a causal effect of experience.
\end{abstract}
\begin{IEEEkeywords}
GitHub Copilot, programmer experience, productivity, secondary analysis, reproducibility
\end{IEEEkeywords}

\section{Introduction}
\label{sec:introduction}
Code-generating assistants introduce an interaction decision alongside the
technical task of producing useful code: assistance can appear automatically or
be requested explicitly. Understanding this distinction requires examining how
programmers use suggestions, rather than assuming that access to a model alone
determines productivity. Barke et al.\ describe interaction patterns in which
programmers use Copilot for both acceleration and exploration~\cite{barke2023grounded}.
Ziegler et al.\ connect perceived productivity with suggestion acceptance,
illustrating that usage measures and productivity judgments need careful
interpretation~\cite{ziegler2022productivity}.

Our starting point is the published study of Copilot interaction modes by
Guglielmi et al.~\cite{guglielmi2026copilot}. We do not claim that the relationship
between experience and suggestion mode is an unexplored topic. Instead, the
question for this replication is whether a transparent participant-level
analysis of the released productivity tables supports an interpretable comparison
of mode-related changes across experience groups. A difference between two
within-group significance decisions would not itself establish a difference
between the groups. We therefore estimate that contrast directly.

We use automatic suggestions as the comparison condition and manual Copilot
invocation as the baseline. This is a comparison of interaction modes, not AI
assistance versus unaided programming. Our exploratory approach combines source
validation, paired productivity changes, and uncertainty estimates. It does not
reproduce every outcome or model in the source study.

\subsection{Contributions and Research Questions}
\label{sec:questions}
The contributions of this course project are:
\begin{itemize}
  \item a reproducible validation pipeline for the published productivity tables,
  including checksum verification and independent metric recalculation;
  \item explicit participant-level estimands for within-group changes and the
  expert-minus-novice contrast, retaining the paired design;
  \item joint reporting of productivity, completion, and time, with documented
  limits on what these measurements can establish.
\end{itemize}
We ask three research questions:
\begin{enumerate}
  \item[RQ1:] What is the mean paired productivity change from manual to automatic
  suggestions among novices?
  \item[RQ2:] What is the corresponding change among experts?
  \item[RQ3:] How does the mean paired change differ between experts and novices?
\end{enumerate}
RQ3 is primary; RQ1 and RQ2 provide context. Section~\ref{sec:methodology}
defines the dataset and estimands. Section~\ref{sec:results} identifies the
remaining results integration, and Section~\ref{sec:limitations} discusses validity.

\section{Methodology}
\label{sec:methodology}
\subsection{Data and Scope}
\label{sec:data}
We analyze the released productivity tables associated with
Guglielmi et al.~\cite{guglielmi2026copilot}. They contain 16 experts and
17 novices, each with one session per mode on different Java maintenance tasks.
We retain the source experience labels and all 66 released observations. The
source investigators collected the data; we recruit no participants and collect
no new human outcomes. The replication-package link, attribution, and SHA-256
checksums are recorded in the repository's data provenance files.

Figure~\ref{fig:pipeline} summarizes our analytical pipeline. Table~\ref{tab:parameters}
lists its implemented parameters. We verify file hashes, normalize decimal-comma
numbers and time-column spellings, require finite values and positive times,
and check completion bounds. We compare recalculated productivity with the
published value using an absolute tolerance of $10^{-8}$. Each participant must
have exactly one observation per mode, consistent group metadata, and different
tasks. Invalid inputs stop processing; they are not silently excluded.

\begin{figure}[t]
\centering
\begingroup\small
\fbox{\begin{minipage}{0.91\columnwidth}
\centering
Published CSV tables and provenance\\
$\downarrow$\\
Checksum, schema, and participant-pair validation\\
$\downarrow$\\
Productivity recalculation and paired changes\\
$\downarrow$\\
Group summaries and stratified bootstrap\\
$\downarrow$\\
CSV summaries and JSON results
\end{minipage}}
\endgroup
\caption{Implemented secondary-analysis pipeline. Metric calculation occurs
after validation. These are participant outcomes, not runtime benchmark logs.}
\label{fig:pipeline}
\end{figure}

\begin{table}[t]
\caption{Implemented analysis parameters in \texttt{configs/config.json}.}
\label{tab:parameters}
\centering
\begin{tabular}{@{}ll@{}}
\toprule
Parameter & Value \\
\midrule
Participants & 33 (17 novice, 16 expert) \\
Observations & 66 sessions \\
Bootstrap repetitions & 10,000 \\
Pseudorandom seed & 42 \\
Confidence level & 95\% \\
Completion-drop flag & Greater than 5 percentage points \\
\bottomrule
\end{tabular}
\end{table}

\subsection{Outcomes and Estimands}
\label{sec:estimands}
Equation~\ref{eq:productivity} defines productivity $P_{im}$ for participant $i$
under mode $m$, where $c_{im}$ is completion in percent and $t_{im}$ is time in
minutes:
\begin{equation}
P_{im}=\frac{c_{im}/100}{t_{im}/60}.
\label{eq:productivity}
\end{equation}
The unit is completed task fraction per hour. Incomplete and zero-completion
sessions remain in the analysis. Session time is not necessarily time to a fully
correct solution. Completion does not measure security or comprehensive code quality.

Equation~\ref{eq:paired} defines the automatic-minus-manual change, and
Equation~\ref{eq:contrast} defines the primary expert-minus-novice contrast:
\begin{equation}
D_i=P_{i,\mathrm{automatic}}-P_{i,\mathrm{manual}}.
\label{eq:paired}
\end{equation}
\begin{equation}
\theta=\bar{D}_{E}-\bar{D}_{N},\qquad
\bar{D}_{g}=\frac{1}{n_g}\sum_{i\in g}D_i.
\label{eq:contrast}
\end{equation}
Here $E$ and $N$ denote expert and novice groups and $n_g$ is their participant
count. Positive $\theta$ indicates a larger automatic-mode benefit for experts.
RQ1 and RQ2 concern $\bar{D}_{N}$ and $\bar{D}_{E}$, respectively; RQ3 concerns
$\theta$. All comparisons are two-sided and exploratory.

We additionally calculate paired completion and time changes. Our student-defined
completion guardrail flags a group mean change below $-5$ percentage points.
It is a reporting rule, not a validated quality threshold. No API error rate or
technical warm-up exclusion is inferred from human session measurements.

\subsection{Uncertainty and Reproducibility}
\label{sec:uncertainty}
The implemented analysis reports group--mode means, medians, and 95th
percentiles. For paired mean changes and $\theta$, it generates 10,000 bootstrap
samples using seed 42. Sampling occurs separately within experience groups and
retains each participant's pair. The interval endpoints are the 2.5th and 97.5th
percentiles of the resulting estimates, with linear interpolation. Sessions are
not resampled as independent participants.

The pipeline uses the Python standard library. The repository documents local
Python 3.12 and a Docker runtime pinned to Python 3.12.12. Running
\texttt{src/benchmark.py} with the full-data configuration generates validated
observations, paired changes, and JSON summaries. These are unadjusted estimates.
Task/order adjustment, leave-one-participant-out checks, the executed analysis
notebook, and final plots remain Task 3 work; they are not presented as completed
methods in this draft.

\section{Results and Discussion}
\label{sec:results}
This section is reserved for the verified Task 3 outputs. The final paper will
map RQ1 and RQ2 to within-group paired changes and RQ3 to the direct contrast
in Equation~\ref{eq:contrast}. A comparison table will report productivity,
completion, and time with dispersion, alongside interval plots. This laboratory
draft intentionally contains no final findings or placeholder numerical results.

\section{Limitations and Conclusion}
\label{sec:limitations}
Experience was not randomly assigned, and the released convenience sample does
not establish population-wide effects. Task and order differences may influence
the unadjusted estimates. Productivity combines completion and time; neither
component establishes comprehensive code quality. Both conditions use Copilot,
so the design cannot estimate the effect of AI access against an AI-free control.
We inspected the published data and findings before specifying this analysis;
our course project is not preregistered. Bootstrap intervals characterize
within-sample uncertainty and do not remove these threats to validity.

The present draft establishes the research questions and reproducible analytical
method. Conclusions about those questions will be written after the Task 3
analysis and sensitivity checks are complete.


\bibliographystyle{IEEEtran}
\bibliography{references}
\end{document}
